\scanpage{72}
\mspage{65}
% archive sheet number: 71
directe est immédiate. De \deleted{fait que} \inserted{l'hyp \deleted{\unsure{inutile|utilisée}} pertinente,} que
\diagram{
\begin{tikzcd}[ampersand replacement=\&,column sep=small,row sep=small]
\&X_1\arrow[dl]\\
X_0\&\\
\&X_2\arrow[ul]
\end{tikzcd}
}
et
\diagram{
\begin{tikzcd}[ampersand replacement=\&,column sep=small,row sep=small]
\&Y_1\arrow[dl]\\
Y_0\&\\
\&Y_2\arrow[ul]
\end{tikzcd}
}
sont fibrants, on n'a besoin d'utiliser qu'une partie \unsure{dont|donc|seult}, viz. que $p_1^X$ et $p_1^Y$ sont fibrants, qui assure\deleted{\unsure{nt|t}} que ces deux carrés $X,Y$ sont $W$-homotopiquement cartésiens, ce qui suffit à impliquer dès que $f_0,f_1,f_2\in W\Rightarrow f_3\in W$.

Je vais essayer de généraliser ce raisonnement, pour traiter le cas d'un ens. ordonné $I$ qui peut être infini, mais tel que $I^\circ$ soit noethérien et sobre, donc toute \deleted{partie} suite croissante de parties fermées de $I$, et toute suite décroissante de parties ouvertes, est stationnaire, et de plus toute partie ouverte filtrante croissante a un plus grand élément.

On veut prouver que $I$ est admissible, et par ``récurrence noethérienne'' on est ramené à prouver ceci: $(\alpha)$ Si $I$ est
% vertical marginal emphasis line beside the condition
\margin{$\Big|$}
tel que pour tout ouvert $I'\subset I$, avec $I'\ne I$, \deleted{\unsure{on}} $I'$ soit admissible, alors $I$ aussi.

\unsure{Ceci|Cela} \deleted{\unsure{ceci}} \inserted{$(\alpha)$} admis, prouvons que $I\inserted{{}=\unsure{I_0|I_\circ}}$ est admissible.
